% Perfil estandar: no activa tagging global.

\documentclass[11pt,a4paper,oneside]{report}
\usepackage[spanish,es-tabla]{babel}
\usepackage[tipografia=serif,interlineado=uno,accesibilidad=standard]{tfl-uba-geologia}
\addbibresource{bibliografia/referencias.bib}
\begin{document}
\chapter{Prueba del estilo bibliografico RAGA}

\section*{Tres formas de citar dentro del texto}

\textbf{1. Cita narrativa.}
\textcite{groeber1946} describió aspectos fundamentales de la geología
regional. Más adelante, \textcite{angelelli1972} analizaron los ciclos
metalogenéticos de la Argentina, mientras que \textcite{groeber1953}
desarrolló la síntesis correspondiente al ámbito ándico.

\textbf{2. Cita parentética.}
El emplazamiento de travertinos puede estar controlado por estructuras
compresivas cuaternarias \parencite{alcacersanchez2020}. Los antecedentes
estratigráficos de la cuenca Neuquina también fueron presentados en reuniones
científicas nacionales \parencite{legarreta1981}.

\textbf{3. Varias fuentes en una misma cita.}
La integración de antecedentes paleontológicos, cartográficos y trabajos
académicos permite construir el marco regional del área de estudio
\parencite{weaver1931,ramos1981,ejemplotesis2025}. La información
institucional consultada en línea se cita de la misma forma
\parencite{inpres2019}.

\clearpage
\printbibliography[title={Referencias de prueba}]
\end{document}
